\documentclass[12pt]{article}
\usepackage[utf8]{inputenc}
\usepackage[T1]{fontenc}
\usepackage[margin=1in]{geometry}
\usepackage{setspace}
\usepackage{xcolor}
\usepackage[colorlinks=true, urlcolor=blue!60!black, linkcolor=blue!60!black]{hyperref}
\setlength{\parindent}{2em}
\setlength{\parskip}{1.25em}
\renewcommand{\baselinestretch}{1.0}
\begin{document}
\begin{center}
{\LARGE \textit{Gonzales v. Raich}, 545 U.S. 1 (2005)\par}
\vspace{0.8em}
{\large Case Brief\par}
\end{center}
\vspace{1em}

\section{Memory Jogger}

Under the Commerce Clause and Necessary and Proper Clause, Congress may criminalize the purely intrastate, noncommercial cultivation and possession of homegrown medical marijuana authorized by state law because regulating local production of a fungible commodity is rationally related to a comprehensive federal regime controlling the interstate drug market (\href{https://www.law.cornell.edu/supct/html/03-1454.ZS.html}{Cornell LII Syllabus}; \href{https://www.courtlistener.com/opinion/799995/gonzales-v-raich/}{CourtListener}).

\section{Detailed Case Facts}

In 1996, California voters approved Proposition 215, codified as the Compassionate Use Act of 1996, Cal. Health \& Safety Code Ann. \S\ 11362.5, making California the first state to authorize limited medicinal use of marijuana (\href{https://www.law.cornell.edu/supct/html/03-1454.ZO.html}{Majority Opinion}). The Act exempted seriously ill patients, their primary caregivers, and recommending physicians from state criminal prosecution for possessing or cultivating marijuana for personal medical purposes upon a physician's recommendation (\href{https://www.law.cornell.edu/supct/html/03-1454.ZO.html}{Majority Opinion}). By the time the case reached the Supreme Court, at least nine states authorized medicinal marijuana use (\href{https://www.law.cornell.edu/supct/html/03-1454.ZO.html}{Majority Opinion}).

Respondents Angel McClary Raich and Diane Monson were California residents suffering from serious medical conditions whose board-certified family physicians concluded, after trying numerous conventional medications, that cannabis was the only effective treatment (\href{https://www.law.cornell.edu/supct/html/03-1454.ZO.html}{Majority Opinion}). Raich's physician attested that forgoing cannabis would cause her excruciating pain and could prove fatal (\href{https://www.law.cornell.edu/supct/html/03-1454.ZO.html}{Majority Opinion}). Monson grew her own marijuana plants at home and ingested the drug by smoking and using a vaporizer; Raich, unable to grow her own, received locally grown cannabis at no charge from two caregivers (litigating as ``John Does''), who also processed it into hashish or keif while Raich processed portions into oils, balms, and foods (\href{https://www.law.cornell.edu/supct/html/03-1454.ZO.html}{Majority Opinion}).

Meanwhile, under Title II of the Comprehensive Drug Abuse Prevention and Control Act of 1970---the federal Controlled Substances Act (CSA), 21 U.S.C. \S\ 801 \textit{et seq.}---Congress established a closed regulatory system classifying marijuana as a Schedule I controlled substance (\S\ 812(c)), making its manufacture, distribution, or possession a federal crime (\S\S\ 841(a)(1), 844(a)) except in FDA-approved research studies (\href{https://www.law.cornell.edu/supct/html/03-1454.ZO.html}{Majority Opinion}). On August 15, 2002, county deputy sheriffs and federal Drug Enforcement Administration (DEA) agents came to Monson's home (\href{https://www.law.cornell.edu/supct/html/03-1454.ZO.html}{Majority Opinion}). Although county officials concluded after a thorough investigation that Monson's use complied completely with California law, DEA agents---following a three-hour standoff---seized and destroyed all six of her cannabis plants (\href{https://www.law.cornell.edu/supct/html/03-1454.ZO.html}{Majority Opinion}).

\section{Procedural History}

Following the DEA raid, Raich, Monson, and Raich's two John Doe caregivers sued the U.S. Attorney General (originally John Ashcroft, later substituted by Alberto R. Gonzales) and the head of the DEA in the U.S. District Court for the Northern District of California (\href{https://www.law.cornell.edu/supct/html/03-1454.ZO.html}{Majority Opinion}; \href{https://www.oyez.org/cases/2004/03-1454}{Oyez}). Seeking declaratory and injunctive relief against CSA enforcement as applied to their personal medical use, respondents alleged violations of the Commerce Clause, the Fifth Amendment Due Process Clause, the Ninth and Tenth Amendments, and the doctrine of medical necessity (\href{https://www.law.cornell.edu/supct/html/03-1454.ZO.html}{Majority Opinion}).

In \textit{Raich v. Ashcroft}, 248 F. Supp. 2d 918 (N.D. Cal. 2003), the District Court denied respondents' motion for a preliminary injunction, finding that despite the severe harm plaintiffs faced, they failed to demonstrate a likelihood of success on the merits (\href{https://www.law.cornell.edu/supct/html/03-1454.ZO.html}{Majority Opinion}). A divided panel of the U.S. Court of Appeals for the Ninth Circuit reversed and ordered entry of a preliminary injunction in \textit{Raich v. Ashcroft}, 352 F.3d 1222 (9th Cir. 2003) (\href{https://www.law.cornell.edu/supct/html/03-1454.ZO.html}{Majority Opinion}). Relying heavily on \textit{United States v. Lopez}, 514 U.S. 549 (1995), and \textit{United States v. Morrison}, 529 U.S. 598 (2000), the Ninth Circuit isolated ``the intrastate, noncommercial cultivation and possession of cannabis for personal medical purposes'' under state law as a distinct, non-economic class of activities beyond Congress's Commerce Clause reach; Judge Beam dissented, viewing the case as indistinguishable from \textit{Wickard v. Filburn}, 317 U.S. 111 (1942) (\href{https://www.law.cornell.edu/supct/html/03-1454.ZO.html}{Majority Opinion}). The Supreme Court granted certiorari on June 28, 2004 (542 U.S. 936), heard oral argument on November 29, 2004, and decided the case on June 6, 2005 (\href{https://www.oyez.org/cases/2004/03-1454}{Oyez}; \href{https://www.law.cornell.edu/supct/html/03-1454.ZS.html}{Cornell LII Syllabus}).

\section{Judicial Votes}

The Supreme Court decided the case by a \textbf{6--3} vote (\href{https://www.law.cornell.edu/supct/html/03-1454.ZS.html}{Cornell LII Syllabus}; \href{https://supreme.justia.com/cases/federal/us/545/1/}{Justia}):
\begin{itemize}
\item \textbf{Majority Opinion (5 Justices):} Justice John Paul Stevens delivered the opinion of the Court, joined by Justices Anthony M. Kennedy, David H. Souter, Ruth Bader Ginsburg, and Stephen G. Breyer (\href{https://www.law.cornell.edu/supct/html/03-1454.ZO.html}{Majority Opinion}).
\item \textbf{Concurring in the Judgment (1 Justice):} Justice Antonin Scalia filed an opinion concurring in the judgment (\href{https://www.law.cornell.edu/supct/html/03-1454.ZC.html}{Scalia Concurrence}).
\item \textbf{Principal Dissent (3 Justices in part):} Justice Sandra Day O'Connor filed a dissenting opinion, in which Chief Justice William H. Rehnquist and Justice Clarence Thomas joined as to all but Part III (\href{https://www.law.cornell.edu/supct/html/03-1454.ZD.html}{O'Connor Dissent}).
\item \textbf{Separate Dissent (1 Justice):} Justice Clarence Thomas filed a lone dissenting opinion (\href{https://www.law.cornell.edu/supct/html/03-1454.ZD1.html}{Thomas Dissent}).
\end{itemize}

\section{Holding}

\textbf{Judgment:} Vacated and remanded (vacating 352 F.3d 1222 and remanding for further proceedings on respondents' unaddressed substantive due process and medical necessity claims) (\href{https://www.law.cornell.edu/supct/html/03-1454.ZS.html}{Cornell LII Syllabus}; \href{https://www.law.cornell.edu/supct/html/03-1454.ZO.html}{Majority Opinion}).

\textbf{Rule Announced:} Congress's authority under the Commerce Clause (U.S. Const. art. I, \S\ 8, cl. 3) and the Necessary and Proper Clause (art. I, \S\ 8, cl. 18) encompasses the power to prohibit the purely local, noncommercial cultivation and possession of a fungible commodity---even for physician-recommended medical use lawful under state law---when Congress has a rational basis to conclude that failing to regulate that intrastate class of activities would undercut a comprehensive federal regulatory scheme governing the interstate market for that commodity (\href{https://www.law.cornell.edu/supct/html/03-1454.ZO.html}{Majority Opinion}).

\textbf{Holding vs. Dicta:} The \textit{holding} is that the CSA's categorical ban on manufacturing and possessing marijuana is constitutional as applied to intrastate, noncommercial medical cultivation and use authorized by California law, because courts may not excise individual applications from a valid comprehensive economic regulatory scheme supported by a rational basis (\href{https://www.law.cornell.edu/supct/html/03-1454.ZO.html}{Majority Opinion}). Key \textit{dicta} includes Justice Stevens's observation that respondents' medical evidence, if credited at trial, would cast serious doubt on Congress's Schedule I finding that marijuana lacks accepted medical use (n. 37), as well as his closing remarks encouraging respondents to pursue administrative rescheduling under 21 U.S.C. \S\ 811 or relief through the democratic process in Congress (\href{https://www.law.cornell.edu/supct/html/03-1454.ZO.html}{Majority Opinion}).

\section{Analysis of Opinions}

\subsection{Majority Opinion (Justice Stevens)}

Justice Stevens framed the case around the third category of Commerce Clause authority recognized in \textit{Perez v. United States}, 402 U.S. 146 (1971), and \textit{Lopez}: Congress's power to regulate intrastate activities that substantially affect interstate commerce (\href{https://www.law.cornell.edu/supct/html/03-1454.ZO.html}{Majority Opinion}). He found the parallels to \textit{Wickard v. Filburn}, 317 U.S. 111 (1942), ``striking'': just as Roscoe Filburn's home-consumed wheat affected supply and demand in a regulated national wheat market, locally grown marijuana is a fungible commodity for which a lucrative interstate market exists (\href{https://www.law.cornell.edu/supct/html/03-1454.ZO.html}{Majority Opinion}). Under deferential rational-basis review, Congress did not need to prove that respondents' conduct in fact substantially affected interstate commerce; given the difficulty of distinguishing locally grown cannabis from interstate cannabis (21 U.S.C. \S\ 801(5)) and the risk of diversion into illicit channels, Congress rationally concluded that exempting homegrown marijuana would leave a ``gaping hole'' in the CSA (\href{https://www.law.cornell.edu/supct/html/03-1454.ZO.html}{Majority Opinion}).

Justice Stevens distinguished \textit{Lopez} and \textit{Morrison} on two grounds. First, those cases involved facial challenges to brief, single-subject statutes regulating noneconomic conduct, whereas respondents mounted an as-applied challenge seeking to excise a slice of a concededly valid comprehensive regulatory regime (\href{https://www.law.cornell.edu/supct/html/03-1454.ZO.html}{Majority Opinion}). Second, citing \textit{Webster's Third New International Dictionary}, Stevens defined ``economics'' as ``the production, distribution, and consumption of commodities,'' concluding that the CSA regulates quintessentially economic activity (\href{https://www.law.cornell.edu/supct/html/03-1454.ZO.html}{Majority Opinion}). Finally, invoking the Supremacy Clause, Stevens rejected the Ninth Circuit's reliance on California's medical-use exemption: state action can neither expand nor circumscribe Congress's plenary commerce power (\href{https://www.law.cornell.edu/supct/html/03-1454.ZO.html}{Majority Opinion}).

\subsection{Concurrence in the Judgment (Justice Scalia)}

Justice Scalia agreed with the outcome but articulated a more nuanced doctrinal foundation rooted explicitly in the Necessary and Proper Clause rather than the Commerce Clause alone (\href{https://www.law.cornell.edu/supct/html/03-1454.ZC.html}{Scalia Concurrence}). Because intrastate activities that merely affect interstate commerce are not themselves interstate commerce, Congress's power to reach them derives from its authority to enact laws ``necessary and proper'' to execute its commerce power (\href{https://www.law.cornell.edu/supct/html/03-1454.ZC.html}{Scalia Concurrence}). Crucially, Scalia argued that under \textit{McCulloch v. Maryland}, 17 U.S. (4 Wheat.) 316 (1819), and \textit{United States v. Darby}, 312 U.S. 100 (1941), Congress may regulate even \textit{noneconomic} local activities---such as simple possession of marijuana---when doing so is an appropriate and ``plainly adapted'' means of making a comprehensive regulation of an interstate market effective (\href{https://www.law.cornell.edu/supct/html/03-1454.ZC.html}{Scalia Concurrence}). Because marijuana is a fungible commodity ``never more than an instant from the interstate market,'' Congress need not rely on California's enforcement controls to prevent diversion, and prohibiting local possession does not violate state sovereignty principles like the anticommandeering rule of \textit{Printz v. United States}, 521 U.S. 898 (1997) (\href{https://www.law.cornell.edu/supct/html/03-1454.ZC.html}{Scalia Concurrence}).

\subsection{Principal Dissent (Justice O'Connor)}

Joined by Chief Justice Rehnquist and Justice Thomas (except as to Part III), Justice O'Connor opened by invoking Justice Brandeis's celebrated vision of states as ``laboratories'' of democracy in \textit{New State Ice Co. v. Liebmann}, 285 U.S. 262 (1932) (\href{https://www.law.cornell.edu/supct/html/03-1454.ZD.html}{O'Connor Dissent}). She argued that the majority's ``comprehensive scheme'' rationale reduced \textit{Lopez} and \textit{Morrison} to ``nothing more than a drafting guide,'' creating a perverse incentive for Congress to bundle questionable local regulations inside sweeping statutes (\href{https://www.law.cornell.edu/supct/html/03-1454.ZD.html}{O'Connor Dissent}). O'Connor criticized the majority's dictionary definition of ``economic'' as ``breathtaking'' and limitless---sweeping in home care, backyard gardening, or charades simply because they substitute for market goods---and distinguished \textit{Wickard} because the Agricultural Adjustment Act exempted small, noncommercial farms and rested on a detailed economic record (\href{https://www.law.cornell.edu/supct/html/03-1454.ZD.html}{O'Connor Dissent}). By contrast, the CSA contained only generalized, unsubstantiated declarations not specific to marijuana, and the government offered no empirical proof that California's tightly regulated medical marijuana patients significantly affected the illicit national market (\href{https://www.law.cornell.edu/supct/html/03-1454.ZD.html}{O'Connor Dissent}). In Part III (writing for herself alone), O'Connor quoted \textit{The Federalist} No. 45 and added that while she would not personally have voted for California's initiative, constitutional federalism required protecting space for state experimentation (\href{https://www.law.cornell.edu/supct/html/03-1454.ZD.html}{O'Connor Dissent}).

\subsection{Separate Dissent (Justice Thomas)}

Justice Thomas wrote separately to advance a thoroughgoing originalist critique of modern Commerce Clause doctrine (\href{https://www.law.cornell.edu/supct/html/03-1454.ZD1.html}{Thomas Dissent}). Citing founding-era dictionaries, \textit{The Federalist}, and scholarship by Randy Barnett, Thomas contended that ``commerce'' originally meant buying, selling, bartering, and transporting goods for trade---excluding productive activities like agriculture or purely personal home cultivation (\href{https://www.law.cornell.edu/supct/html/03-1454.ZD1.html}{Thomas Dissent}). Turning to the Necessary and Proper Clause, Thomas argued that applying the CSA to state-regulated medical marijuana failed both prongs of \textit{McCulloch}: it was neither ``necessary'' (plainly adapted, given California's ID-card registry and anti-diversion rules) nor ``proper'' (because it encroached on states' traditional police powers over criminal law and public health) (\href{https://www.law.cornell.edu/supct/html/03-1454.ZD1.html}{Thomas Dissent}). Mocking the majority's expansive definition of economic activity, Thomas warned that if taken seriously, ``the Federal Government may now regulate quilting bees, clothes drives, and potluck suppers throughout the 50 States,'' rendering the substantial-effects test a ``rootless and malleable standard'' that abandons enumerated powers (\href{https://www.law.cornell.edu/supct/html/03-1454.ZD1.html}{Thomas Dissent}).

\section{Examples: Future Applications}

\begin{enumerate}
\item \textbf{Same-Side (Upheld): Federal Ban on Intrastate Possession of Unregistered Machine Guns.} Congress enacts a comprehensive federal statute regulating the interstate manufacture, sale, and transfer of automatic firearms, including a provision criminalizing purely intrastate, noncommercial assembly and home possession of homemade machine guns. Under \textit{Raich}, this falls squarely within Congress's power because firearms are fungible commodities with an established interstate market, and Congress has a rational basis to conclude that home-assembled weapons could be diverted into interstate commerce and undercut the comprehensive regulatory scheme (\href{https://www.law.cornell.edu/supct/html/03-1454.ZO.html}{Majority Opinion}).

\item \textbf{Same-Side (Upheld): Federal Prohibition on Home Cultivation of State-Legal Psilocybin Mushrooms.} A state decriminalizes the personal home cultivation and therapeutic use of psilocybin mushrooms under a licensed facilitator program. Federal prosecutors charge a resident under the CSA for growing six jars of psilocybin in his basement for personal use. Directly controlled by \textit{Raich}, the CSA's comprehensive regulation of Schedule I controlled substances reaches intrastate noncommercial production of fungible psychoactive substances regardless of state authorization (\href{https://www.law.cornell.edu/supct/html/03-1454.ZO.html}{Majority Opinion}).

\item \textbf{Opposite-Side (Struck Down): Standalone Federal Ban on Intrastate Backyard Composting of Food Scraps.} Congress passes a brief, single-subject statute making it a federal misdemeanor to compost organic kitchen waste within 500 feet of a public school, unaccompanied by any broader federal regulatory scheme governing interstate waste markets or any jurisdictional element. Like the Gun-Free School Zones Act in \textit{Lopez} and the civil remedy in \textit{Morrison}, this standalone regulation targets noncommercial local conduct without being an essential part of a larger interstate economic scheme (\href{https://www.law.cornell.edu/supct/html/03-1454.ZO.html}{Majority Opinion}; \href{https://www.law.cornell.edu/constitution-conan/amendment-10/the-commerce-clause-and-the-tenth-amendment}{Constitution Annotated}).

\item \textbf{Opposite-Side (Struck Down): Federal Mandate Compelling Uninsured Citizens to Purchase Broccoli or Health Insurance.} Congress enacts a comprehensive national nutrition statute that requires every adult resident who does not currently buy fresh vegetables to purchase a minimum monthly quantity of broccoli or pay a penalty. As the Roberts Court later held in \textit{National Federation of Independent Business v. Sebelius}, 567 U.S. 519 (2012), \textit{Raich} extends only to regulating \textit{existing} economic activity (like cultivating a commodity); Congress may not invoke the Commerce Clause or Necessary and Proper Clause to compel individuals to enter commerce by regulating \textit{inactivity} (\href{https://supreme.justia.com/cases/federal/us/567/11-393/case.pdf}{NFIB v. Sebelius}; \href{https://www.law.cornell.edu/wex/enumerated_powers}{Cornell LII Wex}).

\item \textbf{Fence-Sitter (Close Case): Federal Regulation of Intrastate Endangered Species Habitat with No Commercial Value.} Under the Endangered Species Act, federal regulators prohibit a landowner from clearing brush on his private lot to protect a tiny, non-migratory cave beetle that exists solely within one county and has never been bought, sold, or traded. On one hand, under \textit{Raich}'s ``comprehensive regulatory scheme'' and Scalia's Necessary and Proper Clause approach, the government can argue that protecting all listed species is essential to a broader national biodiversity regime affecting interstate commerce in pharmaceuticals and tourism. On the other hand, unlike marijuana or wheat, the cave beetle is not a fungible commodity with an established interstate market, making it far closer to noneconomic local land use outside \textit{Raich}'s core rationale (\href{https://www.law.cornell.edu/supct/html/03-1454.ZO.html}{Majority Opinion}; \href{https://www.law.cornell.edu/supct/html/03-1454.ZC.html}{Scalia Concurrence}).
\end{enumerate}

\section{Critique}

\textbf{Originalist and Textualist Critiques:} Originalist and libertarian constitutional scholars sharply criticized \textit{Raich} as a retreat from the Rehnquist Court's federalism revival in \textit{Lopez} and \textit{Morrison}. Randy E. Barnett---who argued the case for respondents and whose originalist research Justice Thomas cited---contended that \textit{Raich} departed from the founding-era meaning of ``commerce'' as trade and exchange, while noting in subsequent commentary that Scalia's Necessary and Proper Clause concurrence at least preserved textualist constraints against pretextual federal police powers (\href{https://law.lclark.edu/law_reviews/lewis_and_clark_law_review/past_issues/volume_09/number_4.php}{Lewis \& Clark Law Review Symposium}; \href{https://www.law.cornell.edu/supct/html/03-1454.ZD1.html}{Thomas Dissent}). Indeed, the Roberts Court later cabined \textit{Raich} in \textit{NFIB v. Sebelius} (2012) by holding that the commerce power cannot compel commercial activity (\href{https://supreme.justia.com/cases/federal/us/567/11-393/case.pdf}{NFIB v. Sebelius}).

\textbf{Progressive and Doctrinal Critiques:} From a doctrinal and progressive standpoint, \textit{Raich} exposed an ideological irony: the Court's liberal bloc (joined by Kennedy and Scalia) defended expansive New Deal federal power to uphold the War on Drugs, even at the expense of progressive state medical-marijuana reforms (\href{https://ww3.lawschool.cornell.edu/research/JLPP/upload/Somin-507.pdf}{Somin, Cornell J.L. \& Pub. Pol'y}). Commentators in the \textit{Lewis \& Clark Law Review} symposium noted that while \textit{Raich} preserved Congress's capacity to enact national environmental, civil-rights, and health-and-safety programs, it left Commerce Clause doctrine muddled---suggesting that the main operational rule after \textit{Raich} is that local regulations survive precisely because they are embedded in larger statutory schemes (\href{https://law.lclark.edu/law_reviews/lewis_and_clark_law_review/past_issues/volume_09/number_4.php}{Lewis \& Clark Law Review Symposium}; \href{https://loyola-chicago-law-journal.scholasticahq.com/article/76736-the-future-of-the-federalism-revolution-i-gonzales-v-raich-i-and-the-legacy-of-the-rehnquist-court.pdf}{Loyola U. Chi. L.J.}). Scholars also observed that \textit{Raich} prompted some progressives to reconsider the value of federalism as a shield for state-level innovation in healthcare, pain management, and drug policy (\href{https://ww3.lawschool.cornell.edu/research/JLPP/upload/Somin-507.pdf}{Somin, Cornell J.L. \& Pub. Pol'y}).

\textbf{Underlying Competing Values:} At its core, \textit{Raich} pits \textit{national regulatory uniformity and democratic deference to Congress} against \textit{constitutional federalism, enumerated-powers restraint, and state experimentation}. Allowing courts to carve out as-applied exemptions from comprehensive federal market regulations risks undermining uniform national enforcement and inviting ad hoc judicial policymaking; yet deferring categorically to congressional drafting choices risks converting the Commerce and Necessary and Proper Clauses into a plenary federal police power that extinguishes states' traditional authority over local health and criminal law (\href{https://www.law.cornell.edu/supct/html/03-1454.ZO.html}{Majority Opinion}; \href{https://www.law.cornell.edu/supct/html/03-1454.ZD.html}{O'Connor Dissent}).

\section{Key Quotations}

\begin{quote}
``The similarities between this case and \textit{Wickard} are striking. Like the farmer in \textit{Wickard}, respondents are cultivating, for home consumption, a fungible commodity for which there is an established, albeit illegal, interstate market.''
\end{quote}
\noindent --- Justice John Paul Stevens, Opinion of the Court (\href{https://www.law.cornell.edu/supct/html/03-1454.ZO.html}{Cornell LII}; \href{https://www.courtlistener.com/opinion/799995/gonzales-v-raich/}{CourtListener}).

\begin{quote}
``We need not determine whether respondents' activities, taken in the aggregate, substantially affect interstate commerce in fact, but only whether a `rational basis' exists for so concluding.''
\end{quote}
\noindent --- Justice John Paul Stevens, Opinion of the Court (\href{https://www.law.cornell.edu/supct/html/03-1454.ZO.html}{Cornell LII}).

\begin{quote}
``Where necessary to make a regulation of interstate commerce effective, Congress may regulate even those intrastate activities that do not themselves substantially affect interstate commerce.''
\end{quote}
\noindent --- Justice Antonin Scalia, Concurring in the Judgment (\href{https://www.law.cornell.edu/supct/html/03-1454.ZC.html}{Cornell LII}).

\begin{quote}
``If the Court is right, then \textit{Lopez} stands for nothing more than a drafting guide: Congress should have described the relevant crime as `transfer or possession of a firearm anywhere in the nation'---thus including commercial and noncommercial activity, and clearly encompassing some activity with assuredly substantial effect on interstate commerce.''
\end{quote}
\noindent --- Justice Sandra Day O'Connor, Dissenting (\href{https://www.law.cornell.edu/supct/html/03-1454.ZD.html}{Cornell LII}).

\begin{quote}
``Respondents Diane Monson and Angel Raich use marijuana that has never been bought or sold, that has never crossed state lines, and that has had no demonstrable effect on the national market for marijuana. If Congress can regulate this under the Commerce Clause, then it can regulate virtually anything---and the Federal Government is no longer one of limited and enumerated powers.''
\end{quote}
\noindent --- Justice Clarence Thomas, Dissenting (\href{https://www.law.cornell.edu/supct/html/03-1454.ZD1.html}{Cornell LII}).

\end{document}
